
\subsubsection{Permutation Symmetries in Neural Networks}

Neural networks exhibit permutation symmetries arising from exchangeability of computational units (Hecht-Nielsen, 1990; Sussmann, 1992). The Lottery Ticket Hypothesis demonstrated that sparse subnetworks maintain these symmetries during training (Frankle \& Carbin, 2019). Git Re-Basin exploits permutation symmetries through explicit combinatorial optimization to find permutations minimizing distance between model weights (Ainsworth et al., 2022). This approach successfully merges models trained from different initializations but requires explicit search and has only been demonstrated on small models within single architecture families.

\subsubsection{Weight Space Learning and Neural Functional Networks}

Neural Functional Networks (NFN) introduced permutation-equivariant architectures for learning from model zoos (Zhou et al., 2024). NFN uses Deep Sets (Zaheer et al., 2017) to process sets of weights in a permutation-invariant manner, achieving improvements on implicit neural representation datasets. NFN's evaluation focuses on homogeneous model populations—collections with identical architectures. Our work tests whether this approach extends to heterogeneous populations.

\subsubsection{Model Merging and Weight Space Arithmetic}

Task arithmetic (Ilharco et al., 2022) and model merging techniques (Wortsman et al., 2022; Matena \& Raffel, 2022) demonstrate that linearly combining model weights can transfer capabilities or improve performance through operations like θ_merged = αθ₁ + (1-α)θ₂. All successful applications require models to have identical architectures—the weights must align element-wise. Our approach aims to enable merging across architectures by projecting weights into a shared space, but our failure to achieve this suggests per-family canonicalization with post-hoc alignment may be necessary.

\subsubsection{Representation Learning and Equivariance}

Group-equivariant neural networks build equivariance into architectures through group convolutions (Cohen \& Welling, 2016). This explicit construction guarantees equivariance by design. Our MSE equivariance loss attempts to learn permutation invariance through gradient descent. The 0.00\% kernel robustness suggests that MSE-based objectives are insufficient for weight-space permutations. This aligns with the success of explicit group-equivariant architectures and suggests that learned equivariance for weight space may require stronger objectives or explicit group constraints.
